\subsection{Proof theory: natural deduction}
\label{sec:natural-deduction-prop}
\slidesrc{Proof theory (prop.): 1--11}

Truth tables check validity by brute force: they look at \emph{all} valuations. With 20
atomic propositions that is more than a million rows. Humans (and many programs) prove
things differently: they start from what they know and take small, obviously correct
steps until they reach the conclusion. This is what a \textbf{proof system} formalises.

\subsubsection{What is a proof?}
\slidesrc{Proof theory (prop.): 1}

\begin{nbdef}{Proof system, proof, derivation (page 1)}
A \textbf{proof system} is a set of rules for constructing proofs. A \textbf{proof} is a
series of sentences, each of which is either a \textbf{premise} or is justified by applying
one of the rules of the system to \emph{earlier} sentences of the series. A sentence
justified in this way is said to be \textbf{derived} from the premises.
\end{nbdef}

The rules are designed with a very important limitation: they only allow you to derive a
conclusion from premises if the argument is \textbf{deductively valid}, i.e.\ if it is not
possible for all the premises to be true while the conclusion is false (the notion of
logical consequence of Section~\ref{sec:tautologies}). This is what makes it legitimate
to call the series of sentences a proof.

The system used in the course is \textbf{natural deduction}: each rule lets you add a
sentence that must be true given that certain earlier sentences are true, just like you
would convince a friend ``in clear, simple steps''. For each connective there are two kinds
of rules:
\begin{itemize}[nosep]
  \item \textbf{Elim} (elimination) rules break apart a sentence containing the connective
        (from $p \land q$ get $p$);
  \item \textbf{Intro} (introduction) rules build a sentence containing the connective
        (from $p$ and $q$ get $p \land q$);
\end{itemize}
plus a special rule \textbf{Reit} (reiteration: repeat a sentence you already have). These
basic rules are enough to prove \emph{any} truth-functionally valid argument. Since some
proofs become very long, a few \textbf{derived rules} are added later to make proofs shorter
(they prove nothing new).

\begin{nbfact}{How a proof is written (Fitch format)}
Sentences are numbered consecutively. The premises come first, and a horizontal line is
drawn below the last premise. Every following line is a (subsidiary or final) conclusion;
next to it we write the rule that justifies it and the numbers of the earlier lines it is
applied to, e.g.\ ``$\land$\,Elim: 1''. Temporary assumptions open a \textbf{subproof},
drawn further to the right with its own vertical line (see below).
\end{nbfact}

\begin{nbintuition}[Proofs and truth tables are two roads to the same place]
A truth table asks: ``in every possible world where the premises hold, does the
conclusion hold?'' A proof says: ``here is a chain of small steps, each obviously safe,
from the premises to the conclusion''. For propositional logic the two roads lead to
exactly the same valid arguments: the proof system is \textbf{sound} (it never proves an
invalid argument) and \textbf{complete} (every valid argument has a proof).
\end{nbintuition}

\subsubsection{Basic rules without subproofs}
\slidesrc{Proof theory (prop.): 1--2}

\begin{nbfact}{Reit, $\land$\,Elim, $\land$\,Intro, $\limp$\,Elim, $\lor$\,Intro}
\begin{center}
\renewcommand{\arraystretch}{2.6}
\begin{tabularx}{\linewidth}{@{}l c X@{}}
\toprule
\textbf{Rule} & \textbf{Schema} & \textbf{In words} \\
\midrule
Reit & $\dfrac{p}{p}$ & you may repeat any sentence you already have \\
$\land$\,Elim & $\dfrac{p \land q}{p} \quad \dfrac{p \land q}{q}$ & from a conjunction, write either conjunct (as many times as you like) \\
$\land$\,Intro & $\dfrac{p \qquad q}{p \land q}$ & join any two sentences you already have with $\land$ \\
$\limp$\,Elim & $\dfrac{p \limp q \qquad p}{q}$ & from a conditional and its antecedent, write its consequent (\textbf{Modus Ponens}) \\
$\lor$\,Intro & $\dfrac{p}{p \lor q} \quad \dfrac{p}{q \lor p}$ & from any sentence, write a disjunction containing it, with \emph{any} other sentence, in either order \\
\bottomrule
\end{tabularx}
\end{center}
Read a schema as ``if you have the sentences above the line, you may write the one below''.
\end{nbfact}

Why are they safe? $\land$\,Elim: if ``$p$ and $q$'' is true, $p$ is true. $\land$\,Intro:
if both are true, so is their conjunction. $\limp$\,Elim: the only way $p \limp q$ can be
true with $p$ true is that $q$ is true (look at the truth table). $\lor$\,Intro: a
disjunction only needs \emph{one} true disjunct, so adding any other disjunct is harmless.

\begin{nbexample}{Example 1: $\{b \land c,\ j \land k\} \;/\therefore\; c \land k$}
\begin{fitchproof}
\hypo{1}{b \land c} \RPrem
\hypo{2}{j \land k} \RPrem
\have{3}{c} \RAndE{1}
\have{4}{k} \RAndE{2}
\have{5}{c \land k} \RAndI{3,4}
\end{fitchproof}
(Premises are listed in braces; ``$/\therefore$'' reads ``therefore''.)
\end{nbexample}

\begin{nbexample}{Example 2: $\{g,\ h,\ (g \land h) \limp (r \land s)\} \;/\therefore\; s$}
\begin{fitchproof}
\hypo{1}{g} \RPrem
\hypo{2}{h} \RPrem
\hypo{3}{(g \land h) \limp (r \land s)} \RPrem
\have{4}{g \land h} \RAndI{1,2}
\have{5}{r \land s} \RImpE{3,4}
\have{6}{s} \RAndE{5}
\end{fitchproof}
Strategy: the conclusion $s$ is ``inside'' $r \land s$, which is the consequent of line 3.
To unlock it with $\limp$\,Elim we need the antecedent $g \land h$, which we build first.
\end{nbexample}

\begin{nbwarning}[Modus Ponens vs.\ Affirming the Consequent]
$\limp$\,Elim is the valid argument form \textbf{Modus Ponens}: $p \limp q,\ p \;/\therefore\; q$.
A common error is the similar but \textbf{invalid} pattern \textbf{Affirming the
Consequent}: $p \limp q,\ q \;/\therefore\; p$. The truth table of
Section~\ref{sec:tautologies} gives the counterexample ($p$ false, $q$ true).
``If the battery is dead, the car won't start; the car won't start; so the battery is
dead'' is wrong: maybe the tank is empty.
\end{nbwarning}

\begin{nbexample}{Example 3: $\{(a \lor b) \limp c,\ b\} \;/\therefore\; c$}
\begin{fitchproof}
\hypo{1}{(a \lor b) \limp c} \RPrem
\hypo{2}{b} \RPrem
\have{3}{a \lor b} \ROrI{2}
\have{4}{c} \RImpE{1,3}
\end{fitchproof}
From $b$, $\lor$\,Intro could build $a \lor b$ or $b \lor a$ (or $b \lor z$ for any $z$). We
choose $a \lor b$ because that is exactly the antecedent we need for line 1.
\end{nbexample}

\begin{nbexercise}{G1.A}{From English to a proof (page 2--3)}
Identify the premises and conclusion, translate using the letters given, then build a
proof.
\begin{enumerate}[nosep]
  \item The starter must be malfunctioning. The car won't start, but the lights are
        working. If the lights work, the battery must be charged. If the car won't start
        when the battery is charged, there must be a problem with the starter.
        ($s$ = the starter is functioning properly; $c$ = the car starts; $w$ = the lights
        are working; $b$ = the battery is charged.)
  \item If Mia moved out of her parents' house, she must be able to afford to. She can
        afford it only if she found both a roommate and a job. Since she moved out, she
        must have found a job. ($m$ = Mia moved out; $a$ = Mia can afford to move out;
        $r$ = Mia found a roommate; $j$ = Mia found a job.)
\end{enumerate}
\end{nbexercise}
\begin{nbsolution}
\textbf{1.} The conclusion is the first sentence (``the starter must be malfunctioning'' $=
\neg s$); ``but'' is logically an ``and''. Premises: $\neg c \land w$, \ $w \limp b$, \
$(\neg c \land b) \limp \neg s$.
\begin{fitchproof}
\hypo{1}{\neg c \land w} \RPrem
\hypo{2}{w \limp b} \RPrem
\hypo{3}{(\neg c \land b) \limp \neg s} \RPrem
\have{4}{\neg c} \RAndE{1}
\have{5}{w} \RAndE{1}
\have{6}{b} \RImpE{2,5}
\have{7}{\neg c \land b} \RAndI{4,6}
\have{8}{\neg s} \RImpE{3,7}
\end{fitchproof}
\textbf{2.} ``$A$ only if $B$'' means $A \limp B$ (if $A$ happens, $B$ must hold). Premises:
$m \limp a$, \ $a \limp (r \land j)$, \ $m$. Conclusion: $j$.
\begin{fitchproof}
\hypo{1}{m \limp a} \RPrem
\hypo{2}{a \limp (r \land j)} \RPrem
\hypo{3}{m} \RPrem
\have{4}{a} \RImpE{1,3}
\have{5}{r \land j} \RImpE{2,4}
\have{6}{j} \RAndE{5}
\end{fitchproof}
\textbf{Lesson}: ``$A$ only if $B$'' is $A \limp B$, \emph{not} $B \limp A$. ``If'' introduces
the antecedent, ``only if'' introduces the consequent.
\end{nbsolution}

\begin{nbexercise}{G1.B}{Derive the conclusion from the premises (page 3)}
\begin{enumerate}[nosep]
  \item $a,\ a \limp b,\ (a \land b) \limp c \;/\therefore\; b \land c$
  \item $g \land h,\ g \limp (d \limp k),\ h \limp d \;/\therefore\; d \land k$
  \item $a \land b,\ (b \lor c) \limp (d \land e),\ e \limp (b \limp s) \;/\therefore\; s$
  \item $(m \lor o) \limp k,\ k \limp (l \land r),\ m \;/\therefore\; k \land (r \lor s)$
\end{enumerate}
\end{nbexercise}
\begin{nbsolution}
\begin{multicols}{2}
\textbf{1.}
\begin{fitchproof}
\hypo{1}{a} \RPrem
\hypo{2}{a \limp b} \RPrem
\hypo{3}{(a \land b) \limp c} \RPrem
\have{4}{b} \RImpE{2,1}
\have{5}{a \land b} \RAndI{1,4}
\have{6}{c} \RImpE{3,5}
\have{7}{b \land c} \RAndI{4,6}
\end{fitchproof}
\textbf{2.}
\begin{fitchproof}
\hypo{1}{g \land h} \RPrem
\hypo{2}{g \limp (d \limp k)} \RPrem
\hypo{3}{h \limp d} \RPrem
\have{4}{g} \RAndE{1}
\have{5}{h} \RAndE{1}
\have{6}{d \limp k} \RImpE{2,4}
\have{7}{d} \RImpE{3,5}
\have{8}{k} \RImpE{6,7}
\have{9}{d \land k} \RAndI{7,8}
\end{fitchproof}
\columnbreak
\textbf{3.}
\begin{fitchproof}
\hypo{1}{a \land b} \RPrem
\hypo{2}{(b \lor c) \limp (d \land e)} \RPrem
\hypo{3}{e \limp (b \limp s)} \RPrem
\have{4}{b} \RAndE{1}
\have{5}{b \lor c} \ROrI{4}
\have{6}{d \land e} \RImpE{2,5}
\have{7}{e} \RAndE{6}
\have{8}{b \limp s} \RImpE{3,7}
\have{9}{s} \RImpE{8,4}
\end{fitchproof}
\textbf{4.}
\begin{fitchproof}
\hypo{1}{(m \lor o) \limp k} \RPrem
\hypo{2}{k \limp (l \land r)} \RPrem
\hypo{3}{m} \RPrem
\have{4}{m \lor o} \ROrI{3}
\have{5}{k} \RImpE{1,4}
\have{6}{l \land r} \RImpE{2,5}
\have{7}{r} \RAndE{6}
\have{8}{r \lor s} \ROrI{7}
\have{9}{k \land (r \lor s)} \RAndI{5,8}
\end{fitchproof}
\end{multicols}
\textbf{Lesson}: work \emph{backwards} from the conclusion. In 4, $k \land (r \lor s)$ needs
$k$ and $r \lor s$; $r \lor s$ needs only $r$ ($\lor$\,Intro); $r$ is inside $l \land r$;
and so on, until you reach the premises.
\end{nbsolution}

\subsubsection{Rules with subproofs}
\slidesrc{Proof theory (prop.): 3--5}

The remaining Intro and Elim rules need a new ingredient: we temporarily \textbf{assume}
something that was not given as a premise, see what follows, and then give the assumption
up. The part of a proof in which we hold a temporary assumption is a \textbf{subproof},
drawn to the right of a new vertical line (its first line, the assumption, is underlined).

\begin{nbfact}{The subproof restriction (page 10)}
A rule may only be applied to sentences that appear (1) \textbf{earlier} in the proof and
(2) \textbf{outside any subproof that has already ended}. Once a subproof ends, its lines may
depend on an assumption we no longer hold, so they cannot be used individually any more;
the only way to use a finished subproof is as a whole, with one of the four rules below.
Subproofs can be nested, but an inner subproof must end before the outer one.
\end{nbfact}

\begin{nbfact}{$\neg$\,Intro, $\neg$\,Elim, $\limp$\,Intro, $\lor$\,Elim}
\begin{center}
\begin{tabular}{@{}p{6.2cm}@{\hspace{0.8cm}}p{6.2cm}@{}}
\begin{fitchproof}
\open
\hypo{a}{p} \RAss
\have{b}{\vdots}
\have{c}{q}
\have{d}{\neg q}
\close
\have{e}{\neg p} \RNotI{}
\end{fitchproof}
&
\begin{fitchproof}
\open
\hypo{a}{\neg p} \RAss
\have{b}{\vdots}
\have{c}{q}
\have{d}{\neg q}
\close
\have{e}{p} \RNotE{}
\end{fitchproof}
\\[-4pt]
\small $\neg$\,Intro (proof by contradiction) & \small $\neg$\,Elim \\[10pt]
\begin{fitchproof}
\open
\hypo{a}{p} \RAss
\have{b}{\vdots}
\have{c}{q}
\close
\have{e}{p \limp q} \RImpI{}
\end{fitchproof}
&
\begin{fitchproof}
\hypo{z}{p \lor q}
\open
\hypo{a}{p} \RAss
\have{b}{r}
\close
\open
\hypo{c}{q} \RAss
\have{d}{r}
\close
\have{e}{r} \ROrE{}
\end{fitchproof}
\\[-4pt]
\small $\limp$\,Intro (conditional proof) & \small $\lor$\,Elim (proof by cases)
\end{tabular}
\end{center}
\begin{itemize}[nosep]
  \item \textbf{$\neg$\,Intro}: if an assumption leads to a sentence \emph{and} its negation
        (a contradiction), the assumption cannot be true: close the subproof and write its
        negation.
  \item \textbf{$\neg$\,Elim}: exactly like $\neg$\,Intro, but the assumption is a negation
        $\neg p$; from the contradiction, conclude $p$.
  \item \textbf{$\limp$\,Intro}: assume the antecedent $p$; if you derive $q$, close the
        subproof and write $p \limp q$.
  \item \textbf{$\lor$\,Elim}: from $p \lor q$, open two subproofs, one assuming $p$ and one
        assuming $q$; anything $r$ derived in \emph{both} can be written outside them.
\end{itemize}
Citation style: the assumption line through the contradictory pair, e.g.\
``$\neg$\,Intro: 3--(4,5)''; for $\lor$\,Elim, the disjunction and both subproofs,
e.g.\ ``$\lor$\,Elim: 1, 4--6, 7--9''.
\end{nbfact}

\begin{nbintuition}[Everyday versions of these rules]
\begin{itemize}[nosep]
  \item \emph{$\neg$\,Intro, reductio ad absurdum}: ``Suppose the butler did it. Then he was
        in the kitchen at 9. But he was seen in the garden at 9. Contradiction: the butler
        did not do it.''
  \item \emph{$\limp$\,Intro, conditional proof}: ``Suppose I take the 8:00 train. Then I
        arrive at 9:30, so I am in time for the meeting. Hence: if I take the 8:00 train,
        I am in time.'' You never claimed to take the train.
  \item \emph{$\lor$\,Elim, proof by cases}: ``Either it rains or it is sunny. If it rains I
        take the bus; if it is sunny I take the bus (too hot to walk). So I take the bus.''
\end{itemize}
\end{nbintuition}

\begin{nbexample}{Example 4 ($\neg$\,Intro): $\{\neg(g \land h),\ h\} \;/\therefore\; \neg g$}
``It's not the case that both Gary and Harvey will be promoted. Harvey is getting a
promotion. So Gary won't be promoted.''
\begin{fitchproof}
\hypo{1}{\neg(g \land h)} \RPrem
\hypo{2}{h} \RPrem
\open
\hypo{3}{g} \RAss
\have{4}{g \land h} \RAndI{2,3}
\have{5}{\neg(g \land h)} \RReit{1}
\close
\have{6}{\neg g} \RNotI{3--(4,5)}
\end{fitchproof}
Suppose Gary will be promoted (3). Then both will be promoted (4). But we were told they
will not both be promoted (5). So the assumption is false (6). Note the use of Reit: the
rule requires the contradictory pair \emph{inside} the subproof, so we copy line 1 in.
\end{nbexample}

\begin{nbwarning}[Wording in the slides]
In the slides the assumption lines of subproofs are labelled ``Premise''. They are
\emph{temporary} premises (assumptions) that are discharged when the subproof ends, unlike
the real premises at the top. In these notes they are labelled ``Assumption''.
\end{nbwarning}

\begin{nbexample}{Example 5 ($\neg$\,Elim): $\{\neg a \limp n,\ \neg(n \lor e)\} \;/\therefore\; a$}
``If Alicia doesn't get the job, Nina will. But there is no way Nina or Eloise will get the
job. So Alicia will get it.''
\begin{fitchproof}
\hypo{1}{\neg a \limp n} \RPrem
\hypo{2}{\neg(n \lor e)} \RPrem
\open
\hypo{3}{\neg a} \RAss
\have{4}{n} \RImpE{1,3}
\have{5}{n \lor e} \ROrI{4}
\have{6}{\neg(n \lor e)} \RReit{2}
\close
\have{7}{a} \RNotE{3--(5,6)}
\end{fitchproof}
\end{nbexample}

\begin{nbexample}{Example 6 ($\limp$\,Intro): $\{a \limp p,\ p \limp (m \land r),\ m \limp q\} \;/\therefore\; a \limp q$}
``If Alberto gets his way, the policy will be approved. In that case, Marta and Roger will
be angry. Marta will quit if she is angry. So if Alberto gets his way, Marta will quit.''
\begin{fitchproof}
\hypo{1}{a \limp p} \RPrem
\hypo{2}{p \limp (m \land r)} \RPrem
\hypo{3}{m \limp q} \RPrem
\open
\hypo{4}{a} \RAss
\have{5}{p} \RImpE{1,4}
\have{6}{m \land r} \RImpE{2,5}
\have{7}{m} \RAndE{6}
\have{8}{q} \RImpE{3,7}
\close
\have{9}{a \limp q} \RImpI{4--8}
\end{fitchproof}
\end{nbexample}

\begin{nbexample}{Example 7 (nested subproofs): $\{b \limp \neg(c \land d),\ a \limp (b \limp d)\} \;/\therefore\; (a \land b) \limp \neg c$}
\begin{fitchproof}
\hypo{1}{b \limp \neg(c \land d)} \RPrem
\hypo{2}{a \limp (b \limp d)} \RPrem
\open
\hypo{3}{a \land b} \RAss
\have{4}{a} \RAndE{3}
\have{5}{b \limp d} \RImpE{2,4}
\have{6}{b} \RAndE{3}
\have{7}{d} \RImpE{5,6}
\open
\hypo{8}{c} \RAss
\have{9}{c \land d} \RAndI{7,8}
\have{10}{\neg(c \land d)} \RImpE{1,6}
\close
\have{11}{\neg c} \RNotI{8--(9,10)}
\close
\have{12}{(a \land b) \limp \neg c} \RImpI{3--11}
\end{fitchproof}
To prove an implication, assume its antecedent (3) and aim for the consequent $\neg c$; to
prove a negation, assume the un-negated sentence (8) and aim for a contradiction. Inside
the inner subproof we may use lines 1--7, because they are earlier and in subproofs that
are still open.
\end{nbexample}

\begin{nbexample}{Example 8 ($\lor$\,Elim): Rosa arrives late}
``Rosa will get a ride with either Don or Glenn. If Don gives her a ride, she'll be picked up
late and she'll arrive late. If Glenn gives her a ride, he'll make a stop along the way and
Rosa will arrive late. So Rosa will arrive late.''
\begin{fitchproof}
\hypo{1}{d \lor g} \RPrem
\hypo{2}{d \limp (p \land a)} \RPrem
\hypo{3}{g \limp (s \land a)} \RPrem
\open
\hypo{4}{d} \RAss
\have{5}{p \land a} \RImpE{2,4}
\have{6}{a} \RAndE{5}
\close
\open
\hypo{7}{g} \RAss
\have{8}{s \land a} \RImpE{3,7}
\have{9}{a} \RAndE{8}
\close
\have{10}{a} \ROrE{1, 4--6, 7--9}
\end{fitchproof}
Since Rosa will be late either way, we don't need to know \emph{which} alternative of line
1 is true. (The slides contain a typo, ``If Don gives May a ride'': it is Rosa.)
\end{nbexample}

\begin{nbexample}{Example 9 ($\lor$\,Elim): $\{l \limp (h \land u),\ i \limp s,\ l \lor i\} \;/\therefore\; u \lor s$}
``If Luis is given the job, he'll be happy but Igor will be upset. If the job goes to Igor,
Luis will be upset. One of them will get the job, so one of them will be upset.''
\begin{fitchproof}
\hypo{1}{l \limp (h \land u)} \RPrem
\hypo{2}{i \limp s} \RPrem
\hypo{3}{l \lor i} \RPrem
\open
\hypo{4}{l} \RAss
\have{5}{h \land u} \RImpE{1,4}
\have{6}{u} \RAndE{5}
\have{7}{u \lor s} \ROrI{6}
\close
\open
\hypo{8}{i} \RAss
\have{9}{s} \RImpE{2,8}
\have{10}{u \lor s} \ROrI{9}
\close
\have{11}{u \lor s} \ROrE{3, 4--7, 8--10}
\end{fitchproof}
The two cases end with \emph{different} facts ($u$, $s$), so we first make both cases end
with the \emph{same} sentence $u \lor s$ using $\lor$\,Intro.
\end{nbexample}

\begin{nbexercise}{G2}{Derive the conclusion from the premises, basic rules only (page 6)}
\begin{multicols}{2}\small
\begin{enumerate}[nosep]
  \item $\neg b,\ \neg(\neg a \land \neg b) \;/\therefore\; a$
  \item $j \limp (k \land m),\ m \limp \neg k \;/\therefore\; \neg j$
  \item $c \limp (d \land e),\ d \limp \neg h \;/\therefore\; c \limp \neg h$
  \item $(m \lor r) \limp s,\ s \limp (p \limp q) \;/\therefore\; (r \land p) \limp q$
  \item $(b \land c) \limp d,\ (d \land g) \limp h,\ b \land g \;/\therefore\; c \limp h$
  \item $r \limp \neg(p \land q),\ w \limp (r \limp q) \;/\therefore\; w \limp \neg(p \land r)$
  \item $a \lor n,\ n \limp (e \land s),\ a \limp s \;/\therefore\; s$
  \item $j \limp (k \lor l),\ (j \land k) \limp m,\ j,\ l \limp (m \land o) \;/\therefore\; m$
  \item $a \limp (b \lor c),\ d \limp \neg b,\ c \limp \neg(a \land d) \;/\therefore\; a \limp \neg d$
\end{enumerate}
\end{multicols}
\end{nbexercise}
\begin{nbsolution}
All nine arguments are valid (checked by truth table). Strategy hints: to prove $\neg X$,
assume $X$ and find a contradiction ($\neg$\,Intro); to prove $X \limp Y$, assume $X$ and
derive $Y$ ($\limp$\,Intro); to use $X \lor Y$, reason by cases ($\lor$\,Elim).
\begin{multicols}{2}
\textbf{1.}
\begin{fitchproof}
\hypo{1}{\neg b} \RPrem
\hypo{2}{\neg(\neg a \land \neg b)} \RPrem
\open
\hypo{3}{\neg a} \RAss
\have{4}{\neg a \land \neg b} \RAndI{3,1}
\have{5}{\neg(\neg a \land \neg b)} \RReit{2}
\close
\have{6}{a} \RNotE{3--(4,5)}
\end{fitchproof}
\textbf{2.}
\begin{fitchproof}
\hypo{1}{j \limp (k \land m)} \RPrem
\hypo{2}{m \limp \neg k} \RPrem
\open
\hypo{3}{j} \RAss
\have{4}{k \land m} \RImpE{1,3}
\have{5}{k} \RAndE{4}
\have{6}{m} \RAndE{4}
\have{7}{\neg k} \RImpE{2,6}
\close
\have{8}{\neg j} \RNotI{3--(5,7)}
\end{fitchproof}
\textbf{3.}
\begin{fitchproof}
\hypo{1}{c \limp (d \land e)} \RPrem
\hypo{2}{d \limp \neg h} \RPrem
\open
\hypo{3}{c} \RAss
\have{4}{d \land e} \RImpE{1,3}
\have{5}{d} \RAndE{4}
\have{6}{\neg h} \RImpE{2,5}
\close
\have{7}{c \limp \neg h} \RImpI{3--6}
\end{fitchproof}
\columnbreak
\textbf{4.}
\begin{fitchproof}
\hypo{1}{(m \lor r) \limp s} \RPrem
\hypo{2}{s \limp (p \limp q)} \RPrem
\open
\hypo{3}{r \land p} \RAss
\have{4}{r} \RAndE{3}
\have{5}{m \lor r} \ROrI{4}
\have{6}{s} \RImpE{1,5}
\have{7}{p \limp q} \RImpE{2,6}
\have{8}{p} \RAndE{3}
\have{9}{q} \RImpE{7,8}
\close
\have{10}{(r \land p) \limp q} \RImpI{3--9}
\end{fitchproof}
\textbf{5.}
\begin{fitchproof}
\hypo{1}{(b \land c) \limp d} \RPrem
\hypo{2}{(d \land g) \limp h} \RPrem
\hypo{3}{b \land g} \RPrem
\open
\hypo{4}{c} \RAss
\have{5}{b} \RAndE{3}
\have{6}{b \land c} \RAndI{5,4}
\have{7}{d} \RImpE{1,6}
\have{8}{g} \RAndE{3}
\have{9}{d \land g} \RAndI{7,8}
\have{10}{h} \RImpE{2,9}
\close
\have{11}{c \limp h} \RImpI{4--10}
\end{fitchproof}
\end{multicols}
\textbf{6.} Two nested subproofs: $\limp$\,Intro outside, $\neg$\,Intro inside.
\begin{fitchproof}
\hypo{1}{r \limp \neg(p \land q)} \RPrem
\hypo{2}{w \limp (r \limp q)} \RPrem
\open
\hypo{3}{w} \RAss
\have{4}{r \limp q} \RImpE{2,3}
\open
\hypo{5}{p \land r} \RAss
\have{6}{r} \RAndE{5}
\have{7}{q} \RImpE{4,6}
\have{8}{p} \RAndE{5}
\have{9}{p \land q} \RAndI{8,7}
\have{10}{\neg(p \land q)} \RImpE{1,6}
\close
\have{11}{\neg(p \land r)} \RNotI{5--(9,10)}
\close
\have{12}{w \limp \neg(p \land r)} \RImpI{3--11}
\end{fitchproof}
\begin{multicols}{2}
\textbf{7.}
\begin{fitchproof}
\hypo{1}{a \lor n} \RPrem
\hypo{2}{n \limp (e \land s)} \RPrem
\hypo{3}{a \limp s} \RPrem
\open
\hypo{4}{a} \RAss
\have{5}{s} \RImpE{3,4}
\close
\open
\hypo{6}{n} \RAss
\have{7}{e \land s} \RImpE{2,6}
\have{8}{s} \RAndE{7}
\close
\have{9}{s} \ROrE{1, 4--5, 6--8}
\end{fitchproof}
\columnbreak
\textbf{8.}
\begin{fitchproof}
\hypo{1}{j \limp (k \lor l)} \RPrem
\hypo{2}{(j \land k) \limp m} \RPrem
\hypo{3}{j} \RPrem
\hypo{4}{l \limp (m \land o)} \RPrem
\have{5}{k \lor l} \RImpE{1,3}
\open
\hypo{6}{k} \RAss
\have{7}{j \land k} \RAndI{3,6}
\have{8}{m} \RImpE{2,7}
\close
\open
\hypo{9}{l} \RAss
\have{10}{m \land o} \RImpE{4,9}
\have{11}{m} \RAndE{10}
\close
\have{12}{m} \ROrE{5, 6--8, 9--11}
\end{fitchproof}
\end{multicols}
\textbf{9.} The hardest one. Plan: assume $a$, then assume $d$ and look for a contradiction
between $a \land d$ and $\neg(a \land d)$. To get $\neg(a \land d)$ we need $c$, which must
come from $b \lor c$ by cases. In the case $b$ we have both $b$ and $\neg b$, and from a
contradiction anything follows (here $c$, via $\neg$\,Elim).
\begin{fitchproof}
\hypo{1}{a \limp (b \lor c)} \RPrem
\hypo{2}{d \limp \neg b} \RPrem
\hypo{3}{c \limp \neg(a \land d)} \RPrem
\open
\hypo{4}{a} \RAss
\have{5}{b \lor c} \RImpE{1,4}
\open
\hypo{6}{d} \RAss
\have{7}{\neg b} \RImpE{2,6}
\have{8}{a \land d} \RAndI{4,6}
\open
\hypo{9}{b} \RAss
\open
\hypo{10}{\neg c} \RAss
\have{11}{b} \RReit{9}
\have{12}{\neg b} \RReit{7}
\close
\have{13}{c} \RNotE{10--(11,12)}
\close
\open
\hypo{14}{c} \RAss
\have{15}{c} \RReit{14}
\close
\have{16}{c} \ROrE{5, 9--13, 14--15}
\have{17}{\neg(a \land d)} \RImpE{3,16}
\close
\have{18}{\neg d} \RNotI{6--(8,17)}
\close
\have{19}{a \limp \neg d} \RImpI{4--18}
\end{fitchproof}
\textbf{Lesson}: lines 9--16 prove the \emph{disjunctive syllogism} ($b \lor c,\ \neg b
\;/\therefore\; c$) with basic rules only. It is so useful that it becomes a derived rule
in the next sub-section, and with it this proof shrinks from 19 to 11 lines.
\end{nbsolution}

\subsubsection{Derived rules}
\slidesrc{Proof theory (prop.): 6--8, 11}

\textbf{Derived rules} express patterns of argument that we could prove using only the basic
rules. Anything we prove with them, we could prove without them: they are useful not
because of \emph{what} they let us prove, but because they let us prove it more quickly.

\begin{nbfact}{The derived rules}
\begin{center}
\renewcommand{\arraystretch}{2.6}
\begin{tabularx}{\linewidth}{@{}l c X@{}}
\toprule
\textbf{Rule} & \textbf{Schema} & \textbf{In words} \\
\midrule
Modus Tollens & $\dfrac{p \limp q \qquad \neg q}{\neg p}$ & from a conditional and the negation of its consequent, derive the negation of its antecedent \\
Hypothetical syllogism & $\dfrac{p \limp q \qquad q \limp r}{p \limp r}$ & chain conditionals (any number, in any order) \\
Disjunctive syllogism & $\dfrac{p \lor q \quad \neg p}{q} \quad \dfrac{p \lor q \quad \neg q}{p}$ & from a disjunction and the negation of one disjunct, derive the other \\
Double negation ($\neg\neg$) & $\neg\neg p \equiv p$ & drop or add $\neg\neg$, also inside a longer sentence \\
De Morgan's laws & $\begin{array}{c}\neg(p \land q) \equiv \neg p \lor \neg q \\ \neg(p \lor q) \equiv \neg p \land \neg q\end{array}$ & substitute one form for the other, also inside a longer sentence \\
\bottomrule
\end{tabularx}
\end{center}
In proofs these rules are abbreviated \textbf{MT} (Modus Tollens), \textbf{HS} (hypothetical
syllogism), \textbf{DS} (disjunctive syllogism), $\neg\neg$ and \textbf{DeM} (De Morgan).
Double negation and De Morgan are \textbf{substitution rules}: since the two forms are
logically equivalent (same truth table column), one may replace the other even as a
\emph{part} of a bigger sentence. The other rules apply only to whole lines.
\end{nbfact}

\begin{nbexample}{why Modus Tollens is safe: a proof with basic rules}
\notinslides
\begin{fitchproof}
\hypo{1}{p \limp q} \RPrem
\hypo{2}{\neg q} \RPrem
\open
\hypo{3}{p} \RAss
\have{4}{q} \RImpE{1,3}
\have{5}{\neg q} \RReit{2}
\close
\have{6}{\neg p} \RNotI{3--(4,5)}
\end{fitchproof}
Every use of Modus Tollens can be replaced by these 4 extra lines: that is what ``derived''
means.
\end{nbexample}

\begin{nbexample}{Example 10 (Modus Tollens): Raul's newspaper}
``Raul would have picked up his newspaper if he were home. He must not be home, because he
did not pick up his paper.'' $\{h \limp n,\ \neg n\} \;/\therefore\; \neg h$
\begin{fitchproof}
\hypo{1}{h \limp n} \RPrem
\hypo{2}{\neg n} \RPrem
\have{3}{\neg h} \RMT{1,2}
\end{fitchproof}
This may remind you of the invalid form \emph{affirming the consequent}; the negations make
the difference: Modus Tollens is valid.
\end{nbexample}

\begin{nbexample}{Examples 11--13 (hypothetical and disjunctive syllogism)}
\begin{multicols}{2}
\textbf{Ex.~11}
\begin{fitchproof}
\hypo{1}{a \limp b} \RPrem
\hypo{2}{b \limp c} \RPrem
\hypo{3}{c \limp d} \RPrem
\hypo{4}{d \limp e} \RPrem
\have{5}{a \limp e} \RHS{1--4}
\end{fitchproof}
\columnbreak
\textbf{Ex.~12}
\begin{fitchproof}
\hypo{1}{(j \land k) \limp m} \RPrem
\hypo{2}{\neg l} \RPrem
\hypo{3}{k \lor l} \RPrem
\open
\hypo{4}{j} \RAss
\have{5}{k} \RDS{2,3}
\have{6}{j \land k} \RAndI{4,5}
\have{7}{m} \RImpE{1,6}
\close
\have{8}{j \limp m} \RImpI{4--7}
\end{fitchproof}
\end{multicols}
\textbf{Ex.~13}
\begin{fitchproof}
\hypo{1}{(p \land q) \lor (r \land s)} \RPrem
\hypo{2}{p \limp \neg(q \lor u)} \RPrem
\open
\hypo{3}{p \land q} \RAss
\have{4}{p} \RAndE{3}
\have{5}{q} \RAndE{3}
\have{6}{q \lor u} \ROrI{5}
\have{7}{\neg(q \lor u)} \RImpE{2,4}
\close
\have{8}{\neg(p \land q)} \RNotI{3--(6,7)}
\have{9}{r \land s} \RDS{1,8}
\have{10}{s} \RAndE{9}
\end{fitchproof}

In Example 12 the slides list $j$ as a fourth premise, but it is really the assumption of
the subproof (it is discharged by $\limp$\,Intro in line 8).
\end{nbexample}

\begin{nbexample}{Example 14 (double negation): Ned is sick}
``If Ned wasn't sick or on vacation, he'd be at work. He's not at work and not on vacation,
so he must be sick.'' $\{\neg(s \lor o) \limp w,\ \neg w \land \neg o\} \;/\therefore\; s$
\begin{fitchproof}
\hypo{1}{\neg(s \lor o) \limp w} \RPrem
\hypo{2}{\neg w \land \neg o} \RPrem
\have{3}{\neg w} \RAndE{2}
\have{4}{\neg\neg(s \lor o)} \RMT{1,3}
\have{5}{s \lor o} \RDN{4}
\have{6}{\neg o} \RAndE{2}
\have{7}{s} \RDS{5,6}
\end{fitchproof}
Modus Tollens negates the antecedent $\neg(s \lor o)$, producing $\neg\neg(s \lor o)$;
double negation removes the two $\neg$.
\end{nbexample}

\begin{nbexample}{Example 15 (De Morgan): Javier is sick}
``If Javier wasn't sick or on vacation, he'd be at work. He's neither at work nor on
vacation, so he must be sick.'' $\{\neg(s \lor o) \limp w,\ \neg(w \lor o)\} \;/\therefore\; s$
\begin{fitchproof}
\hypo{1}{\neg(s \lor o) \limp w} \RPrem
\hypo{2}{\neg(w \lor o)} \RPrem
\have{3}{\neg w \land \neg o} \RDM{2}
\have{4}{\neg w} \RAndE{3}
\have{5}{\neg\neg(s \lor o)} \RMT{1,4}
\have{6}{s \lor o} \RDN{5}
\have{7}{\neg o} \RAndE{3}
\have{8}{s} \RDS{6,7}
\end{fitchproof}
\end{nbexample}

\begin{nbwarning}[Typos in Examples 14--16 of the slides]
In Example 15 the slides justify lines 4 and 7 with ``$\land$\,Elim: 2'', but line 2 is
$\neg(w \lor o)$, not a conjunction: the conjunction is line 3 (obtained by De Morgan), so
the correct citation is ``$\land$\,Elim: 3'' (as written above). In Example 14 the premise
set is printed as $\{\neg(s\lor o) \supset w,\ \neg w\}$, while the proof uses
$\neg w \land \neg o$ (and $\supset$ is just another symbol for $\limp$). In Example 16,
``Modus Tollens: 4.5'' means lines 4,5.
\end{nbwarning}

\begin{nbexample}{Examples 16--17 (De Morgan inside longer proofs)}
\textbf{Ex.~16}: ``Ed couldn't have passed history. If he had passed history or sociology,
he would have graduated. If he'd graduated, he would have gone to the ceremony, but he
didn't.'' $\{(h \lor s) \limp g,\ (g \limp c) \land \neg c\} \;/\therefore\; \neg h$
\begin{multicols}{2}
\begin{fitchproof}
\hypo{1}{(h \lor s) \limp g} \RPrem
\hypo{2}{(g \limp c) \land \neg c} \RPrem
\have{3}{g \limp c} \RAndE{2}
\have{4}{(h \lor s) \limp c} \RHS{1,3}
\have{5}{\neg c} \RAndE{2}
\have{6}{\neg(h \lor s)} \RMT{4,5}
\have{7}{\neg h \land \neg s} \RDM{6}
\have{8}{\neg h} \RAndE{7}
\end{fitchproof}
\columnbreak
\textbf{Ex.~17}
\begin{fitchproof}
\hypo{1}{a \limp (b \land c)} \RPrem
\hypo{2}{c \limp d} \RPrem
\open
\hypo{3}{\neg d} \RAss
\have{4}{\neg c} \RMT{2,3}
\have{5}{\neg b \lor \neg c} \ROrI{4}
\have{6}{\neg(b \land c)} \RDM{5}
\have{7}{\neg a} \RMT{1,6}
\close
\have{8}{\neg d \limp \neg a} \RImpI{3--7}
\end{fitchproof}
\end{multicols}
\end{nbexample}

\begin{nbexercise}{G3}{Derive the conclusion, derived rules allowed (page 9)}
\begin{multicols}{2}\small
\begin{enumerate}[nosep]
  \item $(a \land b) \limp c,\ \neg c \land \neg d \;/\therefore\; \neg a \lor \neg b$
  \item $\neg(j \land k) \limp m,\ j,\ \neg(l \lor m) \;/\therefore\; \neg k$
  \item $d \limp c,\ c \limp \neg g,\ g \lor h \;/\therefore\; d \limp h$
  \item $j \limp k,\ \neg l \limp \neg(k \lor m) \;/\therefore\; j \limp l$
  \item $\neg b \limp \neg c,\ \neg(a \lor b),\ (\neg a \land \neg c) \limp d \;/\therefore\; d$
  \item $\neg g \limp (h \lor e),\ \neg h \land \neg e,\ g \lor m \;/\therefore\; m$
  \item $\neg(a \lor b) \limp c,\ \neg(b \lor c) \;/\therefore\; \neg a$
  \item $r \limp s,\ g \limp h,\ (s \lor u) \limp g \;/\therefore\; r \limp h$
  \item $o \limp j,\ j \limp (k \lor m),\ \neg(j \land m) \;/\therefore\; o \limp k$
\end{enumerate}
\end{multicols}
\end{nbexercise}
\begin{nbsolution}
\textbf{Watch out: exercises 2, 6 and 7 are not valid as printed in the slides.} A truth
table check finds, for each of them, a valuation that makes all premises true and the
conclusion false, so \emph{no} correct proof can exist:
\begin{itemize}[nosep]
  \item \textbf{2}: $j = k = \true$, $l = m = \false$. The premises actually prove $k$ (not
        $\neg k$): the conclusion was probably meant to be $k$.
  \item \textbf{6}: $g = \true$, $h = e = m = \false$. The premises prove $g$, and $g \lor m$
        then says nothing about $m$. With the third premise $\neg g \lor m$ (``$g \limp m$'')
        the argument becomes valid.
  \item \textbf{7}: $a = \true$, $b = c = \false$. The premises prove $a$ (not $\neg a$).
\end{itemize}
Below, the valid versions are proved.
\begin{multicols}{2}
\textbf{1.}
\begin{fitchproof}
\hypo{1}{(a \land b) \limp c} \RPrem
\hypo{2}{\neg c \land \neg d} \RPrem
\have{3}{\neg c} \RAndE{2}
\have{4}{\neg(a \land b)} \RMT{1,3}
\have{5}{\neg a \lor \neg b} \RDM{4}
\end{fitchproof}
\textbf{2} (conclusion corrected to $k$).
\begin{fitchproof}
\hypo{1}{\neg(j \land k) \limp m} \RPrem
\hypo{2}{j} \RPrem
\hypo{3}{\neg(l \lor m)} \RPrem
\have{4}{\neg l \land \neg m} \RDM{3}
\have{5}{\neg m} \RAndE{4}
\have{6}{\neg\neg(j \land k)} \RMT{1,5}
\have{7}{j \land k} \RDN{6}
\have{8}{k} \RAndE{7}
\end{fitchproof}
\textbf{3.}
\begin{fitchproof}
\hypo{1}{d \limp c} \RPrem
\hypo{2}{c \limp \neg g} \RPrem
\hypo{3}{g \lor h} \RPrem
\open
\hypo{4}{d} \RAss
\have{5}{d \limp \neg g} \RHS{1,2}
\have{6}{\neg g} \RImpE{5,4}
\have{7}{h} \RDS{3,6}
\close
\have{8}{d \limp h} \RImpI{4--7}
\end{fitchproof}
\columnbreak
\textbf{4.}
\begin{fitchproof}
\hypo{1}{j \limp k} \RPrem
\hypo{2}{\neg l \limp \neg(k \lor m)} \RPrem
\open
\hypo{3}{j} \RAss
\have{4}{k} \RImpE{1,3}
\have{5}{k \lor m} \ROrI{4}
\have{6}{\neg\neg(k \lor m)} \RDN{5}
\have{7}{\neg\neg l} \RMT{2,6}
\have{8}{l} \RDN{7}
\close
\have{9}{j \limp l} \RImpI{3--8}
\end{fitchproof}
\textbf{5.}
\begin{fitchproof}
\hypo{1}{\neg b \limp \neg c} \RPrem
\hypo{2}{\neg(a \lor b)} \RPrem
\hypo{3}{(\neg a \land \neg c) \limp d} \RPrem
\have{4}{\neg a \land \neg b} \RDM{2}
\have{5}{\neg a} \RAndE{4}
\have{6}{\neg b} \RAndE{4}
\have{7}{\neg c} \RImpE{1,6}
\have{8}{\neg a \land \neg c} \RAndI{5,7}
\have{9}{d} \RImpE{3,8}
\end{fitchproof}
\end{multicols}
\begin{multicols}{2}
\textbf{6} (third premise corrected to $\neg g \lor m$).
\begin{fitchproof}
\hypo{1}{\neg g \limp (h \lor e)} \RPrem
\hypo{2}{\neg h \land \neg e} \RPrem
\hypo{3}{\neg g \lor m} \RPrem
\have{4}{\neg(h \lor e)} \RDM{2}
\have{5}{\neg\neg g} \RMT{1,4}
\have{6}{m} \RDS{3,5}
\end{fitchproof}
\textbf{7} (conclusion corrected to $a$).
\begin{fitchproof}
\hypo{1}{\neg(a \lor b) \limp c} \RPrem
\hypo{2}{\neg(b \lor c)} \RPrem
\have{3}{\neg b \land \neg c} \RDM{2}
\have{4}{\neg c} \RAndE{3}
\have{5}{\neg\neg(a \lor b)} \RMT{1,4}
\have{6}{a \lor b} \RDN{5}
\have{7}{\neg b} \RAndE{3}
\have{8}{a} \RDS{6,7}
\end{fitchproof}
\columnbreak
\textbf{8.}
\begin{fitchproof}
\hypo{1}{r \limp s} \RPrem
\hypo{2}{g \limp h} \RPrem
\hypo{3}{(s \lor u) \limp g} \RPrem
\open
\hypo{4}{r} \RAss
\have{5}{s} \RImpE{1,4}
\have{6}{s \lor u} \ROrI{5}
\have{7}{g} \RImpE{3,6}
\have{8}{h} \RImpE{2,7}
\close
\have{9}{r \limp h} \RImpI{4--8}
\end{fitchproof}
\textbf{9.}
\begin{fitchproof}
\hypo{1}{o \limp j} \RPrem
\hypo{2}{j \limp (k \lor m)} \RPrem
\hypo{3}{\neg(j \land m)} \RPrem
\open
\hypo{4}{o} \RAss
\have{5}{j} \RImpE{1,4}
\have{6}{k \lor m} \RImpE{2,5}
\have{7}{\neg j \lor \neg m} \RDM{3}
\have{8}{\neg\neg j} \RDN{5}
\have{9}{\neg m} \RDS{7,8}
\have{10}{k} \RDS{6,9}
\close
\have{11}{o \limp k} \RImpI{4--10}
\end{fitchproof}
\end{multicols}
In exercise 6 (line 5), Modus Tollens gives $\neg\neg g$, which is exactly the negation of
the first disjunct $\neg g$ of line 3, so disjunctive syllogism applies directly.
\textbf{Lesson}: before trying to prove something, a quick truth-table check (or a
counterexample) can save you a lot of time; proof rules can only prove valid arguments.
\end{nbsolution}

\begin{nbsummary}[The proof rules at a glance (pages 10--11)]
\small
\begin{tabularx}{\linewidth}{@{}l X@{}}
\textbf{General restriction} & a rule applies only to sentences that are earlier and
outside any subproof that has already ended \\
Reit & repeat any sentence you already have \\
$\land$\,Elim / $\land$\,Intro & take a conjunct / join two sentences with $\land$ \\
$\limp$\,Elim (Modus Ponens) & from $p \limp q$ and $p$, write $q$ \\
$\lor$\,Intro & from $p$, write $p \lor q$ or $q \lor p$ (any $q$) \\
$\neg$\,Intro / $\neg$\,Elim & assumption $p$ (resp.\ $\neg p$) leads to $q$ and $\neg q$:
write $\neg p$ (resp.\ $p$) \\
$\limp$\,Intro & subproof from $p$ to $q$: write $p \limp q$ \\
$\lor$\,Elim & from $p \lor q$ and subproofs $p \leadsto r$, $q \leadsto r$: write $r$ \\
\textit{Derived} & Modus Tollens, hypothetical syllogism, disjunctive syllogism, $\neg\neg$,
De Morgan
\end{tabularx}
\end{nbsummary}
